\begin{abstract}
Et-Tousy, Zyane \& Sharif (2026), ``Adaptive QoS Management in OneM2M Standard: Machine Learning and Deep Learning for IoT Network Optimization'' (Journal of Network and Systems Management), trains a centralized Random Forest on pooled QoS telemetry (RTT, CPU, RAM, success rate) to decide whether IoT traffic should stay local, partially offload, or fully offload to the cloud, reaching 98\% accuracy. The paper explicitly names a gap it does not resolve: its centralized learning mechanism ``might not scale well or guarantee data privacy in distributed IoT environments,'' with federated learning proposed only as future work. This project reproduces the centralized baseline on synthetic QoS telemetry from 6 non-IID simulated IoT edge sites (each dominated by one of the paper's own uniform/burst/real-time traffic scenarios), then builds and evaluates the federated alternative the paper names but does not build: one Random Forest per site, trained on local data only, combined at inference time by averaging predicted probabilities, so no raw telemetry ever leaves its site. Across 5 seeds, the federated ensemble reaches 99.5\% accuracy and 99.5\% macro-F1 versus the centralized baseline's 99.9\%/99.96\% -- a small utility cost for removing the need to centralize raw data -- while a single site's own local model used in isolation, with no federation at all, reaches only 97.6\% accuracy and 97.4\% macro-F1. Federation recovers almost all of the gap between going it alone and pooling everything centrally, without requiring the centralization the baseline paper's own limitations section flags as a privacy risk.
\end{abstract}
